% Figure: dimensionless mass flux through a converging nozzle against the
% back-to-stagnation pressure ratio. The flux rises as the back pressure
% falls, reaches a maximum at the critical pressure ratio (Mach 1 at the
% throat), and then holds flat: the flow is choked and no further drop in
% back pressure raises the throat mass flow. Curve for gamma = 1.3 (steam).
\begin{figure}[htbp]
\centering
\begin{tikzpicture}
\begin{axis}[
  width=12cm, height=7.5cm,
  xlabel={back-pressure ratio $p_b/p_0$},
  ylabel={mass flux $\dot{m}/\dot{m}_{\max}$},
  xmin=0, xmax=1.0, ymin=0, ymax=1.1,
  axis lines=left,
  xtick={0,0.2,0.4,0.546,0.8,1.0},
  xticklabels={$0$,$0.2$,$0.4$,$p^*/p_0$,$0.8$,$1.0$},
  tick label style={font=\scriptsize},
  label style={font=\small},
  legend pos=south west,
  legend style={font=\scriptsize},
  clip=true,
]
% Unchoked branch: rises from (1,0) to the peak at the critical ratio.
% Approximated by a smooth arc peaking at 1.0 at r = 0.546.
\addplot[engblue, very thick, smooth] coordinates {
  (1.0,0) (0.95,0.32) (0.90,0.46) (0.85,0.57) (0.80,0.66) (0.75,0.75)
  (0.70,0.83) (0.65,0.90) (0.60,0.96) (0.546,1.0)
};
\addlegendentry{unchoked (subsonic throat)}
% Choked branch: flat at maximum for back-pressure ratios below the critical.
\addplot[engred, very thick] coordinates {(0.546,1.0) (0,1.0)};
\addlegendentry{choked (sonic throat)}
% critical pressure ratio marker
\addplot[enggray, dashed, thick] coordinates {(0.546,0) (0.546,1.0)};
\node[anchor=south, font=\scriptsize, engred, rotate=90] at (axis cs:0.52,0.35)
  {$p^*/p_0=0.546$};
% choked region label
\node[font=\scriptsize\itshape, engred] at (axis cs:0.25,0.9)
  {choked: flow fixed by throat};
\end{axis}
\end{tikzpicture}
\caption[Choked flow through a converging nozzle]{The mass flux through a
converging nozzle or a relief orifice, normalised by its maximum, plotted
against the back-pressure ratio $p_b/p_0$ for a vapour with $\gamma=1.3$. As
the back pressure falls the flux climbs, reaches its maximum when the throat
reaches the local sound speed at the critical pressure ratio
$p^*/p_0=(2/(\gamma+1))^{\gamma/(\gamma-1)}=0.546$, and then holds flat: once
the throat is sonic, no further drop in back pressure can raise the mass
flow, because pressure signals cannot travel upstream against a sonic throat.
The flat choked branch is why a relief-valve mass rate is set by the sonic
vapour properties at the set pressure and not by the downstream pressure, the
fact the propane relief-sizing case study leaned on.}
\label{fig:vol04ch02:nozzle-choked-flow}
\end{figure}
